\section{Background}
\label{sec:background}

RNG connects $n$ routers directly, with no spine layer, each with $d$ fabric links wired so the graph behaves like a uniformly random $d$-regular graph; at the paper's defaults, which our models use, $n = 1000$ and $d = 64$. Shortest paths on such a graph use only a handful of the $d$ links between two routers, so Spraypoint spreads each pair's traffic in two phases (\autoref{fig:rings}):
\begin{enumerate}
  \item \textbf{Spray.} The source router $s$ hashes the flow to one of its $d$ neighbors, ignoring the destination.
  \item \textbf{Point.} Around each destination $t$, routers form rings: WP0 holds $t$'s $d$ neighbors, WP1 the $p = 4$ waypoints each WP0 router picks among its neighbors, the inner ring (IR) every other router touching WP1, and the outer ring (OR) the rest, almost none. Each router splits traffic over $h = 2$ next hops in the next ring inward, so traffic walks inward one ring per hop.
\end{enumerate}

\begin{figure}[ht]
\centering
\begin{tikzpicture}[scale=0.72, every node/.style={font=\small}]
  \draw[fill=black!4] (0,0) circle (4.2);
  \draw[fill=black!8] (0,0) circle (3.0);
  \draw[fill=black!14] (0,0) circle (1.8);
  \node[circle, fill=black, text=white, inner sep=1.5pt] (t) at (0,0) {$t$};
  \node at (0, 1.35) {WP0};
  \node at (0, 2.45) {WP1};
  \node at (0, 3.65) {IR};
  \node[circle, draw, fill=white, inner sep=1.5pt] (s) at (-6.2, -0.6) {$s$};
  \node[circle, draw, fill=white, inner sep=1pt] (v) at (-3.6, -0.9) {$v$};
  \node[circle, draw, fill=white, inner sep=1pt] (w1) at (-2.4, -0.3) {};
  \node[circle, draw, fill=white, inner sep=1pt] (w0) at (-1.2, -0.2) {};
  \draw[-{Stealth}, thick] (s) -- node[above, font=\scriptsize] {spray} (v);
  \draw[-{Stealth}, thick] (v) -- (w1);
  \draw[-{Stealth}, thick] (w1) -- (w0);
  \draw[-{Stealth}, thick] (w0) -- (t);
  \node[circle, draw, fill=white, inner sep=1pt] (n) at (-1.0, -1.1) {};
  \draw[-{Stealth}, thick, dashed, color=green!50!black] (s) to[bend right=25] node[below, font=\scriptsize] {bypass} (n);
  \draw[-{Stealth}, thick, dashed, color=green!50!black] (n) -- (t);
  \node[align=left, font=\scriptsize, anchor=west] at (4.8, 1.0) {Spraypoint (solid):\\$s \to v \to$ WP1 $\to$ WP0 $\to t$,\\4 hops};
  \node[align=left, font=\scriptsize, color=green!50!black, anchor=west] at (4.8, -1.2) {Bypass (dashed):\\shortest path, here\\$s \to$ WP0 $\to t$, 2 hops};
\end{tikzpicture}
\caption{Spraypoint's rings around a destination $t$. The spray lands on a random neighbor $v$ of the source $s$, usually in the large inner ring, and the flow then walks inward one ring per hop. The bypass (dashed) takes a shortest path instead.}
\label{fig:rings}
\end{figure}

\paragraph{The paper's models.} On a random graph a router's neighbor is close to a uniformly random router, so the spray lands in a ring with probability that ring's size over $n$. We use three of the paper's results (§7 and Appendices E--G of~\cite{rng}).
\begin{itemize}
  \item \textbf{Path length} (Model 7.2). Most paths take 4 hops, because most routers are in IR:
\begin{center}
\small
\begin{tabular}{@{}lccccc|c@{}}
\toprule
Spray lands in & $t$ & WP0 & WP1 & IR & OR & \\
Hops & 1 & 2 & 3 & 4 & 5 & Mean \\
Probability & 0.1\% & 6.4\% & 25.6\% & 67.9\% & $\approx 0$ & \textbf{3.61} \\
\bottomrule
\end{tabular}
\end{center}
  \item \textbf{Path diversity} (Model 7.1). In exchange, each pair gets about $d(1 - e^{-h}) \approx 55$ of a possible 64 edge-disjoint paths, and traffic reaches $t$ through many of its neighbors (high \emph{fan-in}).
  \item \textbf{Throughput} (Model 7.3). The paper measures throughput by \emph{oversubscription} $r$: every router can send $1/r$ of its full rate at once without overloading any link. The worst case is a \emph{permutation}, where every router sends to one partner~\cite{namyar}, and the paper's LP, which splits each pair's traffic over its paths in the best way, gives $r = 3.25$. Model 7.3 approximates the LP by filling the shortest paths first. With $\mu_i d$ the rate, in links' worth, that $s$ can send to $t$ on $i$-hop paths, $r = 1/\sumu$:
\begin{center}
\small
\begin{tabular}{@{}lccccc@{}}
\toprule
& $\mu_2$ & $\mu_3$ & $\mu_4$ & $\mu_5$ & Sum \\
\midrule
Value & 0.064 & 0.131 & 0.111 & $\approx 0$ & \textbf{0.306} \\
Links' worth & 4.1 & 8.4 & 7.1 & 0 & \textbf{19.6} \\
\bottomrule
\end{tabular}
\end{center}
  so $r \approx 3.27$.\footnote{The paper's printed $\mu_5$ formula uses the share of routers in IR where the share in OR is meant; we read it as a typo, and our implementation uses OR.} The $\mu_2$ paths, $s \to \text{WP0} \to t$, about $d^2/n = 4.1$ per pair, are also the pair's shortest paths.
\end{itemize}

The best split is a bound: switches split traffic \emph{evenly} over their next hops (ECMP), so we model the bypass under both (\secref{sec:models}).
